\documentclass[twocolumn,10pt]{article}
\ifdefined\FONTS\else\def\FONTS{ot1}\fi
\def\fontsT{t1}\def\fontsL{lm}\def\fontsP{ptm}
\ifx\FONTS\fontsT \usepackage[T1]{fontenc}\fi
\ifx\FONTS\fontsL \usepackage[T1]{fontenc}\usepackage{lmodern}\fi
\ifx\FONTS\fontsP \usepackage[T1]{fontenc}\usepackage{mathptmx}\usepackage[scaled=0.9]{helvet}\usepackage{courier}\fi
\usepackage[utf8]{inputenc}

\usepackage{amsmath}
\title{Measuring Reading Order in Two Column Layouts}
\author{C. Reader}
\date{}
\begin{document}
\maketitle
\section{Introduction}
Two column layouts are common in scientific papers. A converter must read the left column completely
before moving to the right column, otherwise sentences from different paragraphs are mixed together.
This paper describes a simple test document for that purpose. Each paragraph contains ordinary words so that
any mixing between columns is easy to detect when the extracted text is compared with the source.
\section{Method}
We place several paragraphs of plain text in each column. The paragraphs contain numbers such as 42 and 1999,
hyphenated words such as state-of-the-art, and a short equation $E = mc^2$ to check inline mathematics.
The left column ends with a long paragraph that continues across the column break, which tests whether the
reader joins the two halves of the paragraph in the correct order.
\section{Results}
In our experiments the extracted text matched the source in most cases. Errors were found mainly near
headings and footnotes, where small fonts and short lines confuse the column detection.
Long words such as internationalisation and characterisation were sometimes split by hyphenation at
the end of a line, which a careful reader should join again.
\section{Discussion}
The results suggest that column detection is reliable when the gutter between the columns is wide.
Narrow gutters, figures that span both columns, and footnotes remain difficult. Future versions of the test
will include tables, figures, and footnotes in both columns.
\section{Conclusion}
A small synthetic document is enough to catch most reading order mistakes. We recommend adding such a document
to any regression suite for text extraction.
\end{document}
